\subsubsection{Hidden Liquidity and Dark Trading}

\begin{definition} \hlt{Dark Pool}\\
A dark pool is a trading venue where available liquidity is not displayed before execution. It is mainly used to reduce signalling, market impact, and exposure to opportunistic short-horizon traders.
\end{definition}

\begin{principle} \hlt{Reference Price Trading}\\
Dark pools usually do not form prices independently. Equity dark trades often execute at the midpoint or another reference price from a reliable lit market, so dark liquidity depends on lit price formation.
\end{principle}

\begin{taxonomy} \hlt{Dark Liquidity Pools}
\begin{enumerate}[label=\roman*.]
\setlength{\itemsep}{0pt}
\item \hlt{crossing networks}: match buy and sell orders, often at the primary midpoint;
\item \hlt{internalisation pools}: brokers cross client flow or trade against broker inventory;
\item \hlt{electronic market-makers}: operators or selected liquidity providers accept or reject incoming flow;
\item \hlt{exchange dark books}: hidden order types inside regulated venues;
\item \hlt{dark aggregators}: route one order across several dark pools;
\item \hlt{block pools}: focus on large institutional orders and natural liquidity.
\end{enumerate}
\end{taxonomy}

\begin{mechanism} \hlt{Integrated vs Midpoint Books}\\
Large hidden orders may interact with lit and hidden liquidity in integrated books. Smaller dark orders are typically crossed at the midpoint with basic limit, MAQ, and IOC controls.
\end{mechanism}

\begin{caveat} \hlt{Dark Pools and Price Discovery}\\
Dark pools can reduce market impact, but they rely on lit quotes for reference prices. If too much flow moves dark, public price discovery may weaken; if dark flow is selective and low-impact, lit market quality can improve.
\end{caveat}

\begin{empiricalfinding} \hlt{Dark Activity Drivers}\\
Dark activity is usually concentrated in liquid stocks, high-volume days, volatile periods, and high-spread names. It tends to be lower when order imbalances are large relative to volume.
\end{empiricalfinding}

\begin{mechanism} \hlt{Why Investors Use Dark Pools}\\
The main incentives are secrecy and price improvement. A midpoint execution lets a liquidity taker avoid paying the full spread while hiding the order from traders who might infer and trade against the parent order.
\end{mechanism}

\begin{caveat} \hlt{Dark Liquidity Is Conditional}\\
Dark liquidity is hard to estimate: misses are often invisible, and fills can still be adversely selected.
\end{caveat}

\begin{caveat} \hlt{Toxic Liquidity and Gaming}\\
Dark pools can be gamed through probing, leakage, and adverse selection. Protections include minimum sizes, access controls, and filtering toxic flow.
\end{caveat}

\begin{taxonomy} \hlt{Bilateral vs Multilateral Liquidity}\\
Multilateral books let many participants provide liquidity under common rules. Bilateral trading preserves selection and relationship control, which can help avoid toxic counterparties but may reduce open price discovery.
\end{taxonomy}
